\section{Introduction}\label{sec:introduction}

In his 1982 problem survey, Yau asked whether a complete noncompact
K\"ahler manifold with positive holomorphic bisectional curvature must
be biholomorphic to complex Euclidean space \cite[\S9(b), p.~45]{YauProblems}.
We prove the conjecture under its pointwise positivity hypothesis.

\begin{theorem}\label{thm:uniformization}
Let $M$ be a connected noncompact complex manifold of complex dimension
$n\geq1$. If $M$ admits a smooth complete K\"ahler metric with strictly
positive holomorphic bisectional curvature, then $M$ is biholomorphic
to $\C^n$.
\end{theorem}

Strict positivity in Theorem~\ref{thm:uniformization} is pointwise.
In particular, the theorem assumes neither a uniform positive lower
bound nor a finite upper bound for the curvature.  No volume-growth,
noncollapsing, Ricci-pinching, or topological hypothesis is added.
Its conclusion identifies the entire complex manifold with $\C^n$;
it does not assert that the given metric is Euclidean.
In particular, $M$ is Stein and contractible.
It resolves Yau's uniformization conjecture for strictly positive
holomorphic bisectional curvature in every complex dimension.

\subsection{Antecedents and the remaining analytic difficulties}

The one-dimensional case belongs to the classical relation between
curvature and conformal type. A complete noncompact Riemann surface
of positive Gaussian curvature is conformally the plane, by
Cohn-Vossen's topology theorem and the parabolicity theorem of
Blanc--Fiala--Huber \cite[Satz~8]{CohnVossen}\cite[Theorem~15]{Huber}.
In the compact setting, the Frankel conjecture was resolved by
Mori and by Siu--Yau: positive holomorphic bisectional curvature
characterizes complex projective space \cite{MoriAmple,SiuYauFrankel}.
Yau's noncompact question asks for an analogous characterization
of complex Euclidean space while allowing unrestricted geometry
at infinity \cite[\S9(b), p.~45]{YauProblems}.

One route to this problem studies holomorphic functions and
algebraic embeddings. Greene--Wu showed that positive real sectional
curvature yields a smooth strictly convex exhaustion and hence
Steinness \cite[Theorem~10]{GreeneWuConvex1976}.
Mok obtained an affine-algebraic embedding under positive bisectional
curvature, maximal volume growth, and quadratic scalar-curvature
decay \cite{MokEmbedding1984}.
These results connect curvature to global function theory while
retaining hypotheses or conclusions different from unrestricted
uniformization. Chen--Zhu proved intrinsic growth and average-curvature
estimates under pointwise positive bisectional curvature
\cite{ChenZhu}. In particular, their linear average-scalar-curvature
bound requires no upper curvature bound, but its constant may
depend on the center. They also pursued algebraic compactification
under bounded positive sectional curvature and a finite top Ricci
integral \cite{ChenZhuPositive2009}.
Ni--Tam developed heat deformation of plurisubharmonic functions
under nonnegative bisectional curvature, without a global upper
curvature bound \cite{NiTamStructure}. Their smoothing theorem
keeps an explicit growth condition on the initial function.

A second route uses geometric flow to construct shrinking holomorphic
charts and then assembles those charts by complex dynamics.
Cao's differential Harnack estimates for K\"ahler--Ricci flow
\cite{CaoHarnack} are an analytic antecedent of this approach.
Chau--Tam established uniformization under bounded nonnegative
bisectional curvature and maximal volume growth
\cite{ChauTamComplex2006}. Their later work developed locally
biholomorphic charts, covering and pseudoconvex-domain conclusions,
and soliton-limit methods under additional curvature-decay or
pinching conditions \cite{ChauTamStein,ChauTamSimply}.
The distinction between a covering, a map onto a domain, and a
biholomorphism onto all of $\C^n$ is essential here.
Liu proved uniformization under nonnegative bisectional curvature and
maximal volume growth, without an upper curvature bound, using a
Gromov--Hausdorff approach and retracting holomorphic vector fields
\cite{LiuUniformization}.  Lee--Tam subsequently obtained the same
conclusion under these hypotheses by smoothing the metric and applying
the bounded-curvature theorem \cite{LeeTam}.  Both results retain
maximal volume growth.

Recent work has removed assumptions at infinity in important special
cases.\footnote{The cited versions also record AI assistance.
Datar, Pingali, and Seshadri credit ChatGPT pro~5.2 with the normalization
observation used in Lemma~5.13, as well as feedback and discussions
\cite[Section~1.1 and Acknowledgments]{DatarPingaliSeshadri}.
Their later paper credits ChatGPT~5.5 Pro with the auxiliary function
and induction for the B\'ezout estimates, and ChatGPT~5.0 in thinking
mode with the heat-comparison estimate recalled as Lemma~2.2
\cite[Sections~1--3]{DatarPingaliSeshadriTop}.
Wu reports ChatGPT-6 assistance in developing the proofs and the
author's verification of the arguments
\cite[Acknowledgments]{WuConditional}.}
Datar, Pingali, and Seshadri
\cite[Theorem~1]{DatarPingaliSeshadri} prove that every complete noncompact
K\"ahler surface with positive real sectional curvature is biholomorphic
to $\mathbb C^2$. Their method uses Lipschitz plurisubharmonic weights of
finite Monge--Amp\`ere mass and weighted holomorphic functions.
For positive bisectional curvature, their Theorem~2 gives the same
conclusion under strong Steinness, contractibility, and simple
connectivity at infinity. Here strong Steinness means the existence of
a smooth plurisubharmonic exhaustion with bounded gradient.
Wu \cite[Corollaries~1.5--1.6]{WuConditional} removes contractibility
from this criterion: strong Steinness and simple connectivity at infinity
suffice. Wu also obtains the conclusion under positive bisectional curvature
when the latter topological condition holds and the real sectional
curvatures are nonnegative outside a compact set. These are surface results; positive real sectional
curvature is a stronger hypothesis than positive holomorphic bisectional
curvature.

The finite-mass approach also leads to higher-dimensional results.
Datar, Pingali, and Seshadri \cite[Theorem~1.1]{DatarPingaliSeshadriTop}
prove finiteness of the integral of the top power of the Ricci form on
complete noncompact K\"ahler manifolds with positive real sectional
curvature. Their theorem also applies under positive bisectional
curvature when a smooth strictly plurisubharmonic exhaustion with
uniformly bounded gradient is available; with positive bounded sectional
curvature, they obtain quasiprojectivity
\cite[Theorem~1.2]{DatarPingaliSeshadriTop}.
Dat \cite{DatHessian} develops an $m$-Hessian capacity approach to
finite-Monge--Amp\`ere weights in complex dimensions at least three;
its proposed uniformization route retains proper-map and multiplicity
assumptions.

A complementary direction relates holomorphic growth to the asymptotic
geometry of the K\"ahler--Ricci flow. Shi
\cite[Theorem~1.6]{ShiMinimalDegrees} proves an identity between the
minimal growth degree of holomorphic volume forms and the refined
minimal degrees of holomorphic functions for complete noncompact K\"ahler
manifolds with nonnegative bisectional curvature and nonsplitting
universal cover. Under maximal volume growth, the flow results
and coordinates in \cite[Theorem~1.7 and Corollary~1.12]{ShiMinimalDegrees}
connect these holomorphic degrees with Lyapunov exponents. This gives a
quantitative connection between the function-theoretic and flow
approaches to uniformization in that growth regime.

Two issues must be addressed in the unrestricted problem.
First, the initial metric can have unbounded curvature and arbitrarily
collapsed regions. A complete bounded-curvature flow theorem or a
global tensor maximum principle cannot therefore be used without
first establishing its hypotheses. Second, shrinking charts alone
can identify the manifold with a proper domain in $\C^n$.
To obtain surjectivity, one needs quantitative control of their
coordinate changes even when contraction rates vary substantially.

\subsection{Structure of the proof}

Write $g$ for the initial metric and fix $o\in M$.
We construct a global K\"ahler--Ricci flow $h_t$ by complete Hermitian
approximations on exhausting domains.  Estimates for the metric and
its logarithmic volume loss
$\rho(x,t)=\log(\det g(x)/\det h_t(x))$ pass to the limit
without global curvature control. A spatial weight adapted to the nonnegative form $g-h_t$
then proves scalar comparison before completeness of $h_t$ is known.
This order is important: completeness is a later consequence of
the Harnack inequality.

The curvature argument takes place on the holomorphic cotangent
bundle. We minimize the boundary squared dual norm over holomorphic
discs, with compact boundary walls and a vanishing area penalty.
The canonical symplectic bivector relates the complex Hessian of
the norm to its time derivative. A quantitative deformation lemma
gives a viscosity inequality for the disc infimum; optimizing a
Hermitian ellipsoid in each fiber converts that inequality into
the scalar comparison already proved. The infimum therefore equals
the evolving squared dual norm. This gives joint positivity in space and
logarithmic time, hence the matrix and scalar Harnack inequalities.

The proof uses three time parameters, each for a different estimate:
\begin{center}
\begin{tabular}{lll}
\toprule
Clock & Definition & Role \\
\midrule
Physical time & $t$ & K\"ahler--Ricci flow \\
Logarithmic time & $\tau=\log t$ & Joint positivity and Harnack estimates \\
Volume clock & $s=\rho(o,t)$ & Chart contraction and polynomial degrees \\
\bottomrule
\end{tabular}
\end{center}
The volume clock becomes strictly increasing by the Ricci positivity
proved in Section~\ref{sec:charts}.  Sections~\ref{sec:discs}--\ref{sec:variation}
use $s$ temporarily for the physical time at a disc center; from
Section~\ref{sec:charts} onward it always denotes the volume clock.

The volume clock measures contraction directly.  For the transition
$F_{s,u}$ from a centered unitary chart at clock $s$ to one at $u\ge s$,
\[
 |\det F_{s,u}'(0)|=e^{-(u-s)/2}.
\]
Writing $R(x,t)$ for the scalar curvature of $h_t$, the speed of the
volume clock in logarithmic time is $q(s)=ds/d\tau=tR(o,t)$.
A static logarithmic-norm estimate bounds the greatest singular length
of an $h_t$-unitary local chart, measured in the initial metric $g$,
by a fixed power of its least singular length.  Together with the
determinant loss, this controls contraction in every direction.
Selected time intervals with Ricci pinching then give a contraction
argument for loops.  Simple connectivity allows the local inverse
charts to continue coherently over exhausting compact sets, and a
plurisubharmonic Liouville argument gives one contraction exponent
valid on every fixed compact set.

Finally, we normalize the volume jets of the transition maps and
construct polynomial automorphisms with constant Jacobian.  The
nonlinear shears used to realize normalized jets have determinant
one; the full models retain their contracting linear determinant.
Intervals of nonuniform contraction are allowed: their degree
cost is charged to growth of $q$ in the volume clock.
Inverse iteration of these models makes the corrected inverse charts
converge to an injective holomorphic map $M\to\C^n$.
A subsequential bound for the degree of the forward compositions,
together with smallness on compact subsets of the image and an
inverse-ball inclusion, rules out a finite maximal radius for centered
balls in the image.
This proves that the limiting coordinate map is onto~all~of~$\C^n$.

These mechanisms are proved in the order in which they are used.
Sections~\ref{sec:preliminaries}--\ref{sec:flow} establish the initial
analytic data and scalar comparison. Sections~\ref{sec:discs}--\ref{sec:ellipsoids}
prove joint positivity by discs and ellipsoids.
Section~\ref{sec:charts} obtains completeness, simple connectivity,
and quantitative shrinking charts.
Sections~\ref{sec:dynamics}--\ref{sec:uniformization} construct the
biholomorphism. In particular, the disc-deformation estimate and
the forward-degree estimate are stated with the uniformities
needed at their later applications.

\subsection{Conventions}

We use the real curvature convention in which
$\Rm(a,b,b,a)$ is the sectional-curvature numerator of the plane
spanned by $a,b$. The hypothesis is
\[
 \Rm(u,v,v,u)+\Rm(u,Jv,Jv,u)>0
\]
for every pair of real unit tangent vectors $u,v$ at every point.
By the K\"ahler symmetries, this is equivalent to Griffiths strict
positivity of $T^{1,0}M$ in complex notation. In complex dimension
one the displayed sum equals Gaussian curvature.

We identify a Hermitian tensor $(a_{i\bar j})$ with its real
$(1,1)$-form, suppressing the factor $\sqrt{-1}$ when no confusion
can result. Thus $\dd f$ denotes the complex Hessian of $f$,
$\Delta_h f=\tr_h\dd f$, and
\[
 \partial_t h_t=-\Ric(h_t),\qquad
 \Ric(h)=-\dd\log\det h.
\]
All Laplacian, curvature, and volume conventions in the proof use
this normalization. The dual of a Hermitian metric includes the
transpose dictated by the holomorphic cotangent pairing; matrix
contractions below retain it explicitly when necessary.

The initial complete metric is $g$, a fixed base point is $o$,
and $r(x)=d_g(o,x)$. Constants may change between estimates.
Subscripts indicate dependencies when those dependencies matter.
For Hermitian forms, $A=O(b)h$ means $-Cbh\leq A\leq Cbh$.
The abbreviation psh means plurisubharmonic. A positive constant
extracted from curvature on a specified compact set is never
understood to be uniform over all of $M$.
